% システム設計と技術仕様

本節では、AnimatorSBTフレームワークの技術仕様を提示する。
具体的には、技術スタック（\cref{subsec:stack}）、
スマートコントラクト設計（\cref{subsec:contract}）、
公開テストネット上にデプロイされたプロトタイプ実装（\cref{subsec:prototype}）、
および既存の制作管理クライアントとの統合設計（\cref{subsec:integration}）について述べる。

% ----- 5.1 技術スタック -----
\subsection{技術スタック}
\label{subsec:stack}

提案システムは以下のコンポーネントを使用する。

\begin{itemize}
  \item \textbf{ブロックチェーン：} Polygon PoS~\cite{gangwal2023layer2}（本番デプロイの対象）。Polygonは、
    低いトランザクションコスト（一般的なガス価格で1トランザクションあたり\$0.01未満）、
    Ethereum Virtual Machine（EVM）互換性、および確立されたエコシステムを理由に選定した。
  \item \textbf{スマートコントラクト：} Solidity 0.8.20、OpenZeppelin Contracts v4.x
    （ERC-721、AccessControl、Counters）を使用。
  \item \textbf{メタデータストレージ：} PinataピンニングサービスによるIPFS。
    メタデータJSONファイルはコンテンツアドレス（content-addressed）方式で管理され、
    ピンニングされた後は不変である。
  \item \textbf{発行サービス：} ブロックチェーンとの連携にethers.js v6、
    IPFSアップロードにPinata SDKを使用するNode.jsバックエンド。
  \item \textbf{制作管理クライアント：} 標準的なウェブフレームワークで構築された
    Progressive Web Application（PWA）であり、フリーランスアニメーター向けの
    カット管理、締切追跡、および収入計算機能を提供する。
  \item \textbf{ウォレット管理（本番ターゲット）：} アカウント抽象化（ERC-4337）による
    埋め込みウォレットを使用し、アニメーターがウォレットソフトウェアを
    インストールしたり秘密鍵を管理したりする必要を排除する。
    現行のプロトタイプでは標準的な外部所有アカウント（Externally Owned Account, EOA）を使用しており、
    ERC-4337の統合は本番デプロイ時に計画されている。
\end{itemize}

% ----- 5.2 スマートコントラクト -----
\subsection{スマートコントラクト設計}
\label{subsec:contract}

\texttt{AnimatorSBT}コントラクトは、OpenZeppelinのERC-721実装を拡張し、
2つの主要な変更を加えたものとして設計されている。すなわち、
\textbf{Soulbinding}と\textbf{ロールベースアクセス制御（role-based access control）}である。

\subsubsection{Soulbindingメカニズム}

転送の制限は、ミント、転送、バーンを含むすべてのトークン移動時に呼び出される
\texttt{\_beforeTokenTransfer}フックのオーバーライドにより強制される。

\begin{lstlisting}[basicstyle=\footnotesize\ttfamily, frame=single,
  caption={転送フックによるSoulbindingの強制。},
  label={lst:soulbind},
  morekeywords={function,address,uint256,internal,virtual,override,if,revert,super}]
function _beforeTokenTransfer(
  address from, address to,
  uint256 tokenId, uint256 batchSize
) internal virtual override {
  // Allow mint (from==0), block transfers
  if (from != address(0)
      && to != address(0)) {
    revert SoulboundTransferDisabled();
  }
  super._beforeTokenTransfer(
    from, to, tokenId, batchSize
  );
}
\end{lstlisting}

さらに、\texttt{approve()}および\texttt{setApprovalForAll()}関数は
無条件にリバートするようオーバーライドされており、
承認ベースの転送メカニズムを防止する。
これにより、SBTはいかなるERC-721準拠のインタフェースを通じても転送できないことが保証される。

\subsubsection{なぜERC-5192を採用しないのか}

ERC-5192標準~\cite{daubenschuetz2022erc5192}は、\texttt{locked()}クエリ関数を追加することで、
Soulbound NFTのための最小限のインタフェースを定義している。
本研究がERC-5192を実装しない理由は2つある。
(1)~ERC-5192はイベントを通じて譲渡不能性を\textit{通知する}のみであり、
コントラクトレベルでの強制は行わない---本アプローチは
\texttt{\_beforeTokenTransfer}においてSoulbindingを強制し、
転送を単に非推奨とするのではなく物理的に不可能にする。
(2)~ERC-5192は承認ベースの転送ベクター
（\texttt{approve}、\texttt{setApprovalForAll}）に対処しないが、
本コントラクトはこれらを明示的にブロックする。
将来のバージョンでは、SBT対応ツールとの相互運用性を向上させるために、
本研究の強制メカニズムに加えてERC-5192互換性
（\texttt{locked()}ビュー関数）を追加する可能性がある。

\subsubsection{ロールベースアクセス制御}

コントラクトは、OpenZeppelinの\texttt{AccessControl}を用いて2つのロールを定義する。

\begin{itemize}
  \item \texttt{DEFAULT\_ADMIN\_ROLE}：発行者ロールの付与・取消が可能。
    システム管理者（TOONIQ）が保持する。
  \item \texttt{ISSUER\_ROLE}：新規SBTのミントおよび既存SBTの失効が可能。
    発行サービスのオペレーショナルウォレットに割り当てられる。
\end{itemize}

この分離により、発行サービスの鍵が漏洩した場合でも、
攻撃者はアクセス制御構造自体を変更できないことが保証される。

\subsubsection{コア関数}

コントラクトは以下の外部関数を提供する。

\begin{itemize}
  \item \texttt{mint(address to, string uri)}：指定されたIPFSメタデータURIを持つ
    単一のSBTをミントする。\texttt{SBTMinted}イベントを発行する。
  \item \texttt{mintBatch(address[] recipients, string[] uris)}：プロジェクト完了時に
    複数のアニメーターに証明書を発行する際の効率化のための
    バッチミント関数。
  \item \texttt{revoke(uint256 tokenId)}：SBTをバーンせずに無効化済みとしてマークし、
    オンチェーンの監査証跡を保持する。
    \texttt{SBTRevoked}イベントを発行する。
  \item \texttt{tokensOfOwner(address owner)}：指定アドレスが保有する
    すべてのトークンIDを返す（オフチェーンクエリ用のビュー関数）。
  \item \texttt{isRevoked(uint256 tokenId)}：失効状態を返す。
  \item \texttt{issuedAt(uint256 tokenId)}：発行タイムスタンプを返す。
\end{itemize}

\subsubsection{ガスコスト分析}

\Cref{tab:gas}に、主要な操作のガス消費量を報告する。
これらはSolidityオプティマイザ有効（200回）の条件下で
Hardhatテスト環境において\textbf{実測}されたものである。

\begin{table}[ht]
\centering
\caption{実測されたガス消費量およびPolygon PoS上での推定コスト
（30 gwei、POL = \$0.22、2026年3月時点）。}
\label{tab:gas}
\footnotesize
\setlength{\tabcolsep}{3pt}
\begin{tabular}{@{}lrrr@{}}
\toprule
\textbf{操作} & \textbf{ガス} & \textbf{POL} & \textbf{USD} \\
\midrule
デプロイ          & 1{,}824{,}117  & 0.0547  & \$0.0120  \\
\texttt{mint()}   & 135{,}505      & 0.0041  & \$0.0009  \\
\texttt{mintBatch(10)} & 1{,}039{,}581 & 0.0312 & \$0.0069 \\
\texttt{revoke()} & 50{,}027       & 0.0015  & \$0.0003  \\
\bottomrule
\end{tabular}
\end{table}

このコストでは、単一の貢献証明書のミントにかかる費用は約\textbf{\$0.0009}であり、
個人のフリーランサーにとっても無視できる水準である。
バッチミントにより、トークンあたりのコストは約\$0.0007まで削減される。
コントラクト全体のデプロイコストは約\$0.012である。

% ----- 5.3 プロトタイプ実装 -----
\subsection{プロトタイプ実装}
\label{subsec:prototype}

設計の妥当性を検証するため、3つのコンポーネントから成る動作するプロトタイプを実装した。
公開テストネット上にデプロイされたスマートコントラクト、発行サービス、および検証ポータルである。

\subsubsection{テストネットデプロイ}

\texttt{AnimatorSBT}コントラクトは、Polygon Amoyテストネット
（Chain ID: 80002）上のアドレス
\texttt{0x2cC3\allowbreak{}1B9E\allowbreak{}FBD5\allowbreak{}dDfB\allowbreak{}5456\allowbreak{}EE4c\allowbreak{}76A0\allowbreak{}fcBC\allowbreak{}A5EE\allowbreak{}E117}
（ブロック35{,}058{,}633）にデプロイされた。\footnote{デプロイトランザクションおよび
その後のすべてのインタラクションは、
\url{https://amoy.polygonscan.com/address/0x2cC31B9EFBD5dDfB5456EE4c76A0fcBCA5EEE117}
において公開的に検証可能である。}
35件のユニットテストから成る包括的なテストスイートをHardhatを用いて開発し、
デプロイ、ミント、バッチミント、Soulbinding強制、
失効、ビュー関数、およびアクセス制御をカバーした。
すべてのテストが正常にパスしている。
田中シナリオ（\cref{subsec:flow}）のエンドツーエンドのデモンストレーションが
テストネット上で実行され、
IPFSにホストされたメタデータを持つSBTのミントに成功し、
すべてのオンチェーン状態（残高、トークンURI、失効状態）が検証された。

\subsubsection{発行サービス}

ethers.js~v6およびPinata SDKを使用したNode.jsベースの発行サービスを実装した。
このサービスは、完全な発行ワークフローを実行する。
すなわち、(1)~\cref{tab:metadata}のスキーマに準拠したメタデータJSONの構築、
(2)~\texttt{evidence\_hash}の計算（シリアライズされたattributesのKeccak-256ハッシュ）、
(3)~Pinataを介したIPFSへのメタデータのアップロード、
および(4)~取得したIPFS~URIによる\texttt{AnimatorSBT.mint()}の呼び出しである。

エンドツーエンドのデモンストレーションスクリプトは、
\cref{subsec:flow}で記述した田中シナリオを再現する。
原画アニメーターが第7話で25カットを完成させ、
システムが検証可能なエビデンスハッシュを伴う貢献記録のSBTを発行する。

\subsubsection{検証ポータル}

軽量な検証ポータルを、CDNから読み込んだethers.jsを使用する
静的HTMLページとして実装した。サーバーサイドのインフラストラクチャは不要である。
ポータルはウォレットアドレスを受け取り、コントラクトの\texttt{tokensOfOwner()}
関数に問い合わせ、各SBTのメタデータをIPFSから取得し、
プロジェクト、役割、タスク数、および日付を含む貢献の詳細を表示する。
ポータルはまた、IPFSメタデータからKeccak-256ハッシュを再計算し、
格納された\texttt{evidence\_hash}フィールドと比較することで、
エビデンスハッシュの検証もサポートする。

% ----- 5.4 制作管理クライアントとの統合 -----
\subsection{制作管理クライアントとの統合}
\label{subsec:integration}

AnimatorSBTは、著者らが開発したフリーランスアニメーター向けの
カット管理・収支管理用Progressive Web Applicationである
\textit{Pocket Seishin}\footnote{\url{https://pocket-seishin.vercel.app}にて公開中。}
との統合を想定して設計されている。
本アプリケーションは公開デプロイされており、アクティブなユーザーが存在する。
カットレベルの作業ログはアニメーターのデバイス上にローカル保存されており、
デフォルトでプライバシーが保護される。
SBTの発行はアニメーターの明示的な同意がある場合にのみ行われる。

計画されている統合は、\cref{subsec:flow}で記述した発行フローに従う。
主要な統合ポイントは以下のとおりである。

\begin{enumerate}
  \item \textbf{ウォレットプロビジョニング：}
    新しいアニメーターがPocket Seishinに登録する際、
    アカウント抽象化を用いて埋め込みウォレットが作成される。
    ウォレットアドレスはユーザープロフィールに保存され、
    SBTミントの\texttt{to}パラメータとして使用される。
    アニメーターがウォレットを直接操作する必要はない。
  \item \textbf{作業ログの集約：}
    Pocket Seishinはすでにコア機能の一部として、
    カットレベルのタスクデータ（割り当て、ステータス遷移、完了タイムスタンプ）を
    記録している。
    ユーザーの同意のもと、このローカルに保存されたデータが
    SBTメタデータの構築に使用される。
  \item \textbf{エビデンスハッシュの計算：}
    発行サービスは、認証期間の作業ログエントリをシリアライズしたもののKeccak-256ハッシュを計算する。
    このハッシュはSBTメタデータ（\texttt{evidence\_hash}フィールド）に含まれ、
    後にユーザーのローカルデータと照合して検証できる。
  \item \textbf{証明書の表示：}
    Pocket Seishinのプロフィールページに「SBT証明書」セクションが追加され、
    アニメーターに発行されたトークンがプロジェクト名、
    役割、タスク数、および検証ポータルへのリンクとともに表示される。
\end{enumerate}

アニメーターの視点からは、SBTワークフロー全体が透過的となる。
アニメーターは従来どおりPMCを使用し続けるだけでよく、
プロジェクトのマイルストーン到達時に貢献証明書が自動的に発行される。
ブロックチェーンの知識、ウォレット管理、追加のデータ入力は一切不要である。
